% =============================================================================
% L2PHL_WEIGHTING_TECHNICAL_NOTE.tex
% Project: L2Phl (Listening to the Philippines) - survey weighting technical note
% Author:  Avralt-Od Purevjav
% Style:   AEJ (AEA document class); house tables via avtables.sty (UKR/FIES practice)
% Compile: tectonic L2PHL_WEIGHTING_TECHNICAL_NOTE.tex
%          (or: pdflatex + bibtex 3-pass). Table fragments in tab/ are written by
%          make_tex_tables.py from the Stata diagnostics workbook.
% =============================================================================

\documentclass[AEJ]{AEA}

\usepackage{natbib}
\draftSpacing{1.5}

\usepackage{url}
\usepackage{xcolor}
\usepackage{graphicx}
\usepackage{amsmath}
\usepackage{bm}
\usepackage{enumitem}
\usepackage{avtables}   % house style: booktabs, dcolumn, threeparttable, \sym, \notesize
\usepackage{tabularx}
\newcolumntype{Y}[1]{>{\raggedright\arraybackslash\hsize=#1\hsize}X}
\usepackage{microtype}
\emergencystretch=2em
\newcolumntype{d}[1]{D{.}{.}{#1}}

\makeatletter
\renewcommand{\fnum@table}{\tablename\nobreakspace\thetable}
\renewcommand{\fnum@figure}{\figurename\nobreakspace\thefigure}
\makeatother

\begin{document}

\title{Survey Weighting for Listening to the Philippines:
Construction, Performance, and Use}
\shortTitle{L2Phl Survey Weighting}
\author{Avralt-Od Purevjav\thanks{%
Purevjav: World Bank, Poverty and Equity Global Practice, East Asia and Pacific;
apurevjav@worldbank.org. This note documents the survey weights for the
\textit{Listening to the Philippines} (L2Phl) 2025 baseline and its CATI phone
panel. All diagnostic figures are produced by a self-checking Stata do-file; the
tables are generated from its output. Official Use Only.%
}}
\date{\today}
\pubMonth{July}
\pubYear{2026}
\pubVolume{}
\pubIssue{}

\begin{abstract}
This note documents how the L2Phl survey weights are constructed, how they
perform, and which weight to apply to each indicator. It covers both data
collection modes: the CAPI face-to-face baseline of September--October 2025, for
which the weights are built once and frozen, and the CATI phone panel, for which
the weights are re-calibrated against census margins each round. The note states
the sampling design and the weighting arithmetic, reports weight dispersion, the
design effect from weighting, effective sample sizes, and reconciliation to
census totals, analyses panel attrition and its effect on representativeness, and
closes with usage guidance and the corresponding survey-estimation commands.
\end{abstract}

\maketitle

% =============================================================================
\section{Overview and Purpose}
% =============================================================================

This note documents how the L2Phl survey weights are constructed, how they
perform, and which weight to apply to each indicator. It covers both data
collection modes: the CAPI face-to-face baseline of September--October 2025, for
which the weights are built once and frozen, and the CATI phone panel, for which
the weights are re-calibrated each round against census margins.

The intended readers are the fieldwork firm, who carry the weighting into
production, and World Bank colleagues, who use the weighted estimates in analysis
and reporting. The note states the design and the arithmetic plainly, so that any
weighted figure can be traced back to its construction and reproduced.

The survey delivers three analytical weights, and each answers a different
question about the same data (Table~\ref{tab:threeweights}). The population weight
$\mathrm{popw}$ is the World Bank convention for welfare and SDG reporting,
because larger households, which are often poorer, carry weight in proportion to
the people who live in them.

\begin{table}[!t]
\raggedright
\begin{threeparttable}
\caption{The three analytical weights}
\label{tab:threeweights}
\begin{tabularx}{\textwidth}{@{}Y{0.6} Y{1.0} Y{1.4}@{}}
\toprule
Weight & Unit of analysis & Reads as \\
\midrule
$\mathrm{indw}$ & Individual (person) & ``X.X\% of individuals\ldots'' or a per-person mean \\
$\mathrm{hhw}$  & Household & ``X.X\% of households\ldots'' or a per-household mean \\
$\mathrm{popw}$ & Population (via household) & ``X.X\% of Filipinos live in households where\ldots'' \\
\bottomrule
\end{tabularx}
\end{threeparttable}
\end{table}

% =============================================================================
\section{CAPI Baseline Sampling Design}
% =============================================================================

The baseline uses a two-stage stratified cluster design. Strata are defined by
region crossed with urban/rural residence, so most regions contribute an urban
and a rural stratum. Five special geographic areas enter as a single stratum
each, so the design comprises 39 strata in total.

At the first stage, primary sampling units---barangays or enumeration
areas---are drawn with probability proportional to size (PPS), where the measure
of size is the 2020 census population of the unit. The frame is the 2020 census
PSU list with official population counts, PSGC identifiers, and urban/rural flags;
units missing a population count or a residence status are excluded, and stratum
totals are computed from the frame. In total, 247 PSUs are selected across strata.

At the second stage, field teams select a fixed number of households per
PSU---ten completed interviews---using a random-walk protocol, with no advance
listing of households. Interviewers start from a random point and follow a fixed
route and skip rule; when an address is ineligible or does not respond, the team
continues under the same rule until ten completed households are reached. This
approximates equal-probability selection of households within the PSU, and any
small departures are absorbed by the measure-of-size correction and
post-stratification described below.

% =============================================================================
\section{CAPI Weight Construction}
% =============================================================================

The base household weight is the inverse of the joint probability that a
household is selected. Within a stratum $s$, the first-stage probability that PSU
$p$ enters the sample is the standard PPS inclusion probability, and at the second
stage households are selected with equal probability inside the PSU:
%
\begin{equation}
\pi^{1}_{sp} = \frac{n_s\,\mathrm{POP}_{sp}}{\mathrm{POP}_s},
\qquad
\pi^{2}_{sp} = \frac{m_{sp}}{\mathrm{HH}_{sp}},
\end{equation}
%
where $n_s$ is the number of sampled PSUs in stratum $s$, $\mathrm{POP}_{sp}$ and
$\mathrm{HH}_{sp}$ the population and household count of PSU $p$, $\mathrm{POP}_s$
the stratum population, and $m_{sp}=10$ the completed interviews per PSU. The
joint inclusion probability is the product, and the base household weight is its
inverse:
%
\begin{equation}
\pi_{sph} = \pi^{1}_{sp}\,\pi^{2}_{sp}
= \frac{n_s\, m_{sp}}{\mathrm{POP}_s}\,\overline{hs}_{sp},
\qquad
w^{\mathrm{base}}_{sph} = \frac{1}{\pi_{sph}}
= \frac{\mathrm{POP}_s}{n_s\, m_{sp}\,\overline{hs}_{sp}},
\end{equation}
%
with $\overline{hs}_{sp}=\mathrm{POP}_{sp}/\mathrm{HH}_{sp}$ the average household
size of PSU $p$. Within a PSU all sampled households share the same base weight.

Because the first stage uses population rather than the number of households as
the measure of size, the base design carries a small mismatch, and the correction
for it is exact. Multiplying by the ratio of the PSU's average household size to
the stratum's average household size converts the PPS-by-population weight to the
weight a PPS-by-households design would have produced:
%
\begin{equation}
w^{\mathrm{adj}}_{sph}
= w^{\mathrm{base}}_{sph}\,\frac{\overline{hs}_{sp}}{\overline{hs}_{s}}
= \frac{\mathrm{HH}_s}{n_s\, m_{sp}},
\qquad
\overline{hs}_{s}=\frac{\mathrm{POP}_s}{\mathrm{HH}_s}.
\end{equation}
%
A PSU with larger-than-average households is scaled up, because PPS-by-population
under-represents it in household terms, while a PSU with smaller households is
scaled down. Each stratum is then aligned to its official census household total
by a single ratio, so that the sum of household weights in every stratum equals
the census count while within-PSU constancy is preserved:
%
\begin{equation}
\mathrm{hhw}_{sph}
= w^{\mathrm{adj}}_{sph}\,
\frac{\mathrm{HH}^{\mathrm{cen}}_s}{\sum_{p,h} w^{\mathrm{adj}}_{sph}}.
\end{equation}

Because every member of a sampled household is in scope, each person initially
carries the household weight. Ratio calibration is then applied on the person file,
cell by cell, so that the weighted count in each age-sex cell $c$ within a stratum
equals its census total; the population weight is the within-household sum of
person weights:
%
\begin{equation}
\mathrm{indw}_i
= \mathrm{hhw}_{sph}\,
\frac{\mathrm{POP}^{\mathrm{cen}}_{s,c}}{\sum_{i\in(s,c)}\mathrm{hhw}_{sph}},
\qquad
\mathrm{popw}_h = \sum_{i\in h}\mathrm{indw}_i.
\end{equation}
%
These steps produce the three outputs. The individual weight $\mathrm{indw}$ is
the age-sex-calibrated person weight; the household weight $\mathrm{hhw}$ is
carried one per household; and the population weight $\mathrm{popw}$ is the
within-household sum of person weights, so that a household-level indicator
weighted by $\mathrm{popw}$ reports the number of people rather than the number of
households. By construction, a person indicator weighted by $\mathrm{indw}$ and a
household indicator weighted by $\mathrm{popw}$ give consistent population totals.

% =============================================================================
\section{CATI Panel Re-weighting}
% =============================================================================

Every CATI round is re-weighted from the frozen CAPI baseline weights rather than
from scratch. The baseline individual weight is the starting point each round, and
the re-weighting re-anchors it to census margins for that round's achieved sample.
The individual weight is rescaled cell by cell so that each census
age-by-sex-by-stratum cell $c$ sums to its census population for round $r$; the
population weight is the within-household sum and the household weight the
within-household mean:
%
\begin{align}
\mathrm{indw}_{i,r}
&= w^{\mathrm{base}}_i\,
\frac{N^{\mathrm{cen}}_{c}}{\sum_{j\in c,\,r} w^{\mathrm{base}}_j},
\\[2pt]
\mathrm{popw}_{h,r} = \sum_{i\in h}\mathrm{indw}_{i,r},
\qquad
&\tilde{h}_{h,r} = \frac{1}{n_h}\sum_{i\in h}\mathrm{indw}_{i,r}
= \frac{\mathrm{popw}_{h,r}}{n_h}.
\end{align}
%
The household weight is then smoothed to the PSU average to damp variation and
finally rescaled so that each stratum $s$ matches its census household total:
%
\begin{equation}
\bar{h}_{h,r} = \frac{1}{|P(h)|}\sum_{h'\in P(h)}\tilde{h}_{h',r},
\qquad
\mathrm{hhw}_{h,r} = \bar{h}_{h,r}\,
\frac{H^{\mathrm{cen}}_{s}}{\sum_{h'\in s}\bar{h}_{h',r}},
\end{equation}
%
where $w^{\mathrm{base}}_i$ is the frozen baseline individual weight,
$N^{\mathrm{cen}}_{c}$ the census population of cell $c$, $n_h$ the size of
household $h$, $P(h)$ its PSU, and $H^{\mathrm{cen}}_{s}$ the census household count
of stratum $s$. Age is grouped into 0--17, 18--45, and 46 and older, crossed with
sex and stratum. Re-anchoring to census margins each round is what absorbs panel
attrition and non-response: a round with fewer completed interviews in a cell
simply carries larger weights there, and each round remains nationally
representative on the age, sex, and geographic margins.

% =============================================================================
\section{Weight Performance and Diagnostics}
% =============================================================================

This section reports how the weights behave. For the CAPI baseline and for each
CATI round it documents weight dispersion, the design effect from weighting
($\mathrm{DEFF}_w = 1 + \mathrm{CV}^2$, where CV is the coefficient of variation of
the weights), the effective sample size ($\mathrm{ESS} = N / \mathrm{DEFF}_w$,
cross-checked against Kish's effective sample size), and the reconciliation of
weighted totals to census benchmarks.

Table~\ref{tab:capidist} reports the CAPI baseline. The household weight is the
most efficient, with a design effect near $1.07$ and an effective sample close to
its nominal size; the population weight, which spreads across household sizes, is
the least. All 39 strata reconcile to their census population and household totals:
the largest absolute residual is $6.3\times 10^{-8}$ for the population total and
$2.3\times 10^{-9}$ for the household total, so the weighted totals match the
census benchmarks to within rounding. Calibration slightly raises dispersion
relative to the design weights (Table~\ref{tab:caleffect}).

\begin{table}[!t]
\raggedright
\begin{threeparttable}
\caption{CAPI baseline: weight distribution and design effect}
\label{tab:capidist}
\footnotesize
\setlength{\tabcolsep}{4pt}
\begin{tabular}{@{}l r r r r r r r r r r@{}}
\toprule
Weight & $N$ & Mean & Median & Min & Max & CV & Max/min & $\mathrm{DEFF}_w$ & ESS & Kish ESS \\
\midrule
indw & 10,496 & 10,353 & 10,515 & 3,271 & 37,264 & 0.35 & 11.39 & 1.12 & 9,374 & 9,374 \\
popw & 2,470 & 43,995 & 41,146 & 3,432 & 206,870 & 0.58 & 60.28 & 1.33 & 1,856 & 1,856 \\
hhw & 2,470 & 10,686 & 11,489 & 4,632 & 24,114 & 0.26 & 5.21 & 1.07 & 2,308 & 2,308
\\
\bottomrule
\end{tabular}
\begin{tablenotes}\notesize
\item \textit{Notes:} $\mathrm{DEFF}_w = 1+\mathrm{CV}^2$; $\mathrm{ESS}=N/\mathrm{DEFF}_w$.
$\mathrm{indw}$ over all persons; $\mathrm{popw}$ and $\mathrm{hhw}$ over households.
\end{tablenotes}
\end{threeparttable}
\end{table}

\begin{table}[!t]
\raggedright
\begin{threeparttable}
\caption{CAPI baseline: effect of calibration (design vs.\ final weights)}
\label{tab:caleffect}
\footnotesize
\begin{tabular}{@{}l l r r r@{}}
\toprule
Weight & Stage & CV & $\mathrm{DEFF}_w$ & ESS \\
\midrule
indw & design & 0.26 & 1.07 & 9,850 \\
indw & final & 0.35 & 1.12 & 9,374 \\
popw & design & 0.54 & 1.29 & 1,912 \\
popw & final & 0.58 & 1.33 & 1,856 \\
hhw & design & 0.26 & 1.07 & 2,311 \\
hhw & final & 0.26 & 1.07 & 2,308
\\
\bottomrule
\end{tabular}
\end{threeparttable}
\end{table}

Table~\ref{tab:catiindw} and Table~\ref{tab:catihhw} report the CATI panel by
round. The design effect stays stable across the eight rounds, near $1.1$--$1.2$
for the household weight, so the per-round re-calibration does not inflate weight
variation as the panel rolls forward.

\begin{table}[!t]
\raggedright
\begin{threeparttable}
\caption{CATI panel: individual weight by round}
\label{tab:catiindw}
\footnotesize
\begin{tabular}{@{}c r r r r r r@{}}
\toprule
Round & $N$ & CV & $\mathrm{DEFF}_w$ & ESS & Kish ESS & $\sum \mathrm{indw}$ \\
\midrule
1 & 5,378 & 0.44 & 1.19 & 4,513 & 4,513 & 108,667,043 \\
2 & 5,353 & 0.39 & 1.15 & 4,654 & 4,654 & 108,667,043 \\
3 & 5,225 & 0.50 & 1.25 & 4,170 & 4,170 & 108,667,043 \\
4 & 5,572 & 0.43 & 1.19 & 4,698 & 4,698 & 108,667,042 \\
5 & 5,347 & 0.45 & 1.20 & 4,447 & 4,447 & 108,667,043 \\
6 & 5,261 & 0.43 & 1.19 & 4,434 & 4,434 & 108,667,043 \\
7 & 5,046 & 0.45 & 1.20 & 4,199 & 4,199 & 108,667,043 \\
8 & 5,013 & 0.49 & 1.24 & 4,047 & 4,047 & 108,667,043
\\
\bottomrule
\end{tabular}
\end{threeparttable}
\end{table}

\begin{table}[!t]
\raggedright
\begin{threeparttable}
\caption{CATI panel: household and population weight by round}
\label{tab:catihhw}
\footnotesize
\begin{tabular}{@{}c r r r r r r@{}}
\toprule
Round & $N_{hh}$ & $\mathrm{CV}_{hhw}$ & $\mathrm{DEFF}_{hhw}$ & $\mathrm{ESS}_{hhw}$ & $\mathrm{CV}_{popw}$ & $\mathrm{ESS}_{popw}$ \\
\midrule
1 & 1,239 & 0.35 & 1.12 & 1,107 & 0.63 & 886.3 \\
2 & 1,193 & 0.31 & 1.09 & 1,091 & 0.57 & 902.5 \\
3 & 1,174 & 0.43 & 1.19 & 988.8 & 0.67 & 812.8 \\
4 & 1,243 & 0.36 & 1.13 & 1,099 & 0.61 & 903.3 \\
5 & 1,184 & 0.36 & 1.13 & 1,047 & 0.61 & 861.9 \\
6 & 1,157 & 0.35 & 1.12 & 1,029 & 0.60 & 848.9 \\
7 & 1,112 & 0.35 & 1.13 & 988.3 & 0.60 & 819.6 \\
8 & 1,115 & 0.36 & 1.13 & 988.0 & 0.61 & 810.6
\\
\bottomrule
\end{tabular}
\end{threeparttable}
\end{table}

\begin{table}[!t]
\raggedright
\begin{threeparttable}
\caption{CATI panel: weighted composition by round}
\label{tab:caticomp}
\footnotesize
\begin{tabular}{@{}c r r r r@{}}
\toprule
Round & $N_{hh}$ & Urban share & $Q1$ share & $Q5$ share \\
\midrule
1 & 1,239 & 0.55 & 0.22 & 0.18 \\
2 & 1,193 & 0.55 & 0.22 & 0.18 \\
3 & 1,174 & 0.55 & 0.22 & 0.18 \\
4 & 1,243 & 0.55 & 0.23 & 0.18 \\
5 & 1,184 & 0.55 & 0.22 & 0.19 \\
6 & 1,157 & 0.55 & 0.23 & 0.18 \\
7 & 1,112 & 0.55 & 0.23 & 0.18 \\
8 & 1,115 & 0.55 & 0.22 & 0.20
\\
\bottomrule
\end{tabular}
\begin{tablenotes}\notesize
\item \textit{Notes:} Urban is a calibration margin, so its weighted share is
held fixed by construction. $Q1$/$Q5$ are the bottom and top income quintiles.
\end{tablenotes}
\end{threeparttable}
\end{table}

% =============================================================================
\section{Panel Attrition and Representativeness}
% =============================================================================

The CATI panel is not a closed cohort. Of the 1{,}239 households reached in Round
1, 705 (56.9 percent) are still interviewed in Round 8, while 534 have left; each
round also brings households first reached later, 410 of them by Round 8
(Table~\ref{tab:retention}). Because the weights re-anchor the achieved sample to
census margins every round, attrition on the calibrated dimensions---age, sex,
stratum, and the household total---is corrected by construction. The open question
is whether the households that leave differ systematically from those that stay on
characteristics the calibration does not target.

\begin{table}[!t]
\raggedright
\begin{threeparttable}
\caption{Retention by round}
\label{tab:retention}
\footnotesize
\begin{tabular}{@{}c r r r r@{}}
\toprule
Round & Present & From R1 cohort & \% R1 retained & New since R1 \\
\midrule
1 & 1,239 & 1,239 & 100 & 0 \\
2 & 1,193 & 860 & 69.41 & 333 \\
3 & 1,174 & 794 & 64.08 & 380 \\
4 & 1,243 & 801 & 64.65 & 442 \\
5 & 1,184 & 767 & 61.90 & 417 \\
6 & 1,157 & 734 & 59.24 & 423 \\
7 & 1,112 & 694 & 56.01 & 418 \\
8 & 1,115 & 705 & 56.90 & 410
\\
\bottomrule
\end{tabular}
\end{threeparttable}
\end{table}

To test for selection, households present in Round 1 are split into those retained
through Round 8 and those that left, and compared on their Round-1 characteristics
(Table~\ref{tab:attrbias}). On the characteristics that matter most for welfare
measurement the two groups are statistically indistinguishable: income quintile,
food insecurity, the head's sex, urban residence, and geographic composition show
no significant stayer-leaver difference (all $p>0.05$). The one significant
difference is household size: retained households are larger by about a third of a
member ($4.48$ against $4.15$, $p=0.005$). Because household size enters the
age-sex calibration indirectly, since larger households contribute more
person-cells, this residual is partly absorbed by the person-weight step. Baseline
employment is not tested here, because the phone panel's employment module has
thin Round-1 coverage.

\begin{table}[!t]
\raggedright
\begin{threeparttable}
\caption{Attrition: stayers versus leavers (Round-1 characteristics)}
\label{tab:attrbias}
\footnotesize
\setlength{\tabcolsep}{5pt}
\begin{tabular}{@{}l r r r c r r c l@{}}
\toprule
Characteristic & Stayer & Leaver & Diff. & Test & Stat. & $p$ & Sig. & Verdict \\
\midrule
urban & 60 & 54.68 & 5.32 & $\chi^2$ & 3.52 & 0.061 & * & similar \\
female-headed & 40 & 40.07 & -0.07 & $\chi^2$ & 0.00 & 0.979 & ns & similar \\
food insecure (mod.+sev.) & 40.14 & 42.70 & -2.55 & $\chi^2$ & 0.82 & 0.366 & ns & similar \\
household size & 4.48 & 4.15 & 0.33 & $t$ & -2.79 & 0.005 & *** & differs \\
income quintile & 2.98 & 3.04 & -0.05 & $t$ & 0.66 & 0.509 & ns & similar \\
geography (stratum) & -- & -- &  & $\chi^2$ & 52.54 & 0.059 & * & similar \\
\textit{overall} &  &  &  &  &  &  &  & BIASED
\\
\bottomrule
\end{tabular}
\begin{tablenotes}\notesize
\item \textit{Notes:} Percentages except household size and income quintile
(means). Geographic composition is tested across the 39 strata (region proxy).
Significance: *** $p<0.01$, ** $p<0.05$, * $p<0.10$.
\end{tablenotes}
\end{threeparttable}
\end{table}

A formal joint test therefore flags the panel as biased, but the selection is
confined to household size and is modest; on income and food security, the
dimensions the reweighting does not calibrate to, the households that remain are
not a selected subset. A direct check on the headline food-security series supports
this. Re-estimating the moderate-or-severe rate on the balanced panel, households
present in all eight rounds, tracks the full-sample series closely (Round 1, $42.4$
against $41.0$ percent; Round 8, $20.4$ against $20.8$), so the recovery in food
security is a within-household improvement rather than an artefact of who leaves
the panel.

% =============================================================================
\section{Usage Guidance: Which Weight, When}
% =============================================================================

The weight to apply follows from the unit of analysis and the framing of the
indicator. Estimates should carry design-based standard errors, so the survey
design is declared once and the \texttt{svy:} prefix is used throughout:
%
\begin{quote}\ttfamily
svyset psu [pweight=hhw], strata(stratum)\\
svyset psu [pweight=indw], strata(stratum)
\end{quote}
%
for household-level and person-level estimates respectively. Population-framed
welfare indicators use \texttt{[pweight=popw]} on the household file. The weight
mapping for the CAPI modules (Table~\ref{tab:usecapi}) and the CATI modules
(Table~\ref{tab:usecati}) follows the unit of analysis and, for welfare framings,
the population weight.

\begin{table}[!t]
\raggedright
\begin{threeparttable}
\caption{CAPI modules: which weight}
\label{tab:usecapi}
\footnotesize
\begin{tabularx}{\textwidth}{@{}Y{0.7} Y{0.7} Y{1.6}@{}}
\toprule
Module & Weight & Basis \\
\midrule
M01 Roster & indw & individual characteristics \\
M02 Education & indw & individual attendance/attainment/spend \\
M03 Employment & indw & individual labour-market outcomes \\
M04 Income & indw / hhw / popw & earnings = indw; HH transfers = hhw; welfare coverage (pension, 4Ps) = popw \\
M05 Finance & hhw (popw for emergency capacity) & HH financial decisions; vulnerability framed per population \\
M06 Migration & hhw / indw & HH intent = hhw; individual displacement = indw \\
M07 Health & indw (popw for hospital bill) & individual care-seeking/coverage; HH expense per population \\
M08 Food & hhw / indw & HH sourcing = hhw; individual SSB consumption = indw \\
M09 Hazards & popw & share of affected population warned/assisted \\
M10 Dwelling & popw & housing conditions (people) \\
M11 Sanitation & popw (hhw for waste behaviour) & SDG-6 people with improved sanitation \\
M12 Utilities & popw & SDG-6/7 people with electricity/water/internet \\
M13 Assets & hhw (popw for cooking fuel) & HH asset stock; clean-cooking exposure per population \\
M14 Views & hhw & HH head's subjective assessment \\
\bottomrule
\end{tabularx}
\end{threeparttable}
\end{table}

\begin{table}[!t]
\raggedright
\begin{threeparttable}
\caption{CATI modules: which weight}
\label{tab:usecati}
\footnotesize
\begin{tabularx}{\textwidth}{@{}Y{0.7} Y{0.7} Y{1.6}@{}}
\toprule
Module & Weight & Basis \\
\midrule
M00 Passport & --- & identifiers/metadata, not analysed \\
M01 Roster & indw & individual characteristics \\
M02 Education & indw & individual attendance \\
M03 Shocks & hhw & household exposure \\
M04 Employment & indw & individual labour-market outcomes \\
M05 Income & hhw (indw for individual earnings) & household transfers; individual earnings \\
M06 Finance & hhw & household financial decisions \\
M07 Health & indw & individual care-seeking/coverage \\
M08 Food (FIES) & hhw & household food security \\
M09 Views & hhw & HH respondent's assessment \\
\bottomrule
\end{tabularx}
\end{threeparttable}
\end{table}

% =============================================================================
\section{Limitations and Caveats}
% =============================================================================

The two modes do not always measure the same thing. The CAPI baseline records
financial usage---money received into or paid out of a bank or mobile-money
account---while the CATI panel records account ownership. Ownership exceeds recent
usage, so the two figures are not comparable, and a rise from the baseline usage
figure to a panel ownership figure is a definitional artefact rather than adoption.
The valid comparison is within the panel, round to round.

Attitudinal items shift with the mode. Life satisfaction and similar subjective
measures run lower on the phone panel than face-to-face. The story is in the panel
band across rounds, because the gap between the CAPI level and the CATI level is a
mode difference rather than a change in welfare.

Small subgroup cells should be read with caution. Cells below about 30
observations are thin, and in the CATI panel some modules run only in selected
rounds, which thins the base further. Disaggregated estimates in those cells carry
wide sampling error.

Headline estimates should report design-based standard errors under the
\texttt{svy:} prefix and cite the effective sample size, so that the precision
behind each figure is transparent.

\end{document}
